% SPDX-License-Identifier: CC-BY-4.0
\documentclass[11pt,letterpaper]{article}

\usepackage[T1]{fontenc}
\usepackage{amsmath,amssymb}
\usepackage{tikz}
\usepackage{pgfplots}
\pgfplotsset{compat=1.18}
\usepackage{algorithm}
\usepackage{algpseudocode}
\usepackage{booktabs}
\usepackage{enumitem}
\usepackage[hidelinks]{hyperref}

\newcommand{\N}{\mathbb{N}}
\newcommand{\Z}{\mathbb{Z}}
\newcommand{\Q}{\mathbb{Q}}
\newcommand{\R}{\mathbb{R}}
\newcommand{\C}{\mathbb{C}}
\newcommand{\Map}[3]{#1 \colon #2 \to #3}
\newcommand{\Mod}[1]{#1\text{-}\mathrm{Mod}}
\newcommand{\GroupHom}{\mathrm{Hom}_{\mathbf{Grp}}}
\newcommand{\RingHom}{\mathrm{Hom}_{\mathbf{Ring}}}

\newtheorem{theorem}{Theorem}
\newtheorem{lemma}[theorem]{Lemma}
\newtheorem{corollary}[theorem]{Corollary}
\newtheorem{definition}[theorem]{Definition}
\newtheorem{remark}[theorem]{Remark}
\newenvironment{proof}{\noindent\textit{Proof.}}{\hfill\(\square\)}

\title{Generalized Collatz Library}
\author{ascendnoosphere}
\date{\today}

\begin{document}
\maketitle

\section*{Return cores of the multiplier \(d-1\) map}

Throughout, fix integers \(d\ge2\) and \(c\ge1\) and define the multiplier
\(q=d-1\) map
\[
 C(n)=\begin{cases} n/d, & d\mid n,\\ (d-1)n+c, & d\nmid n.\end{cases}
\]
Division has precedence, \(C^{k}\) denotes \(k\) original instructions, so
\(C^{0}\) is the identity and \(C^{k+1}=C\circ C^{k}\). A state \(n\) is
\emph{accessible} if \(d\mid n\) or \(n\equiv c\pmod d\), and the state after
the first division is the \emph{return} state
\[
 R(n)=\begin{cases} n/d, & d\mid n,\\ \bigl((d-1)n+c\bigr)/d, & d\nmid n.\end{cases}
\]
The map \(R\) is total, but only at accessible states does one or two
instructions of \(C\) realize it, as Theorem~\ref{Thm:Realize} states. The
\emph{core} is the interval \(\{1,\dots,c\}\). All states are natural numbers,
and a state is \emph{positive} when it is nonzero.

\section{The residues of one instruction}

\begin{lemma}[Affine branch]\label{Lem:AffineBranch}
For every \(n\),
\[
 C(n)=\begin{cases} n/d, & d\mid n,\\ (d-1)n+c, & d\nmid n,\end{cases}
 \qquad
 R(n)=\begin{cases} n/d, & d\mid n,\\ \bigl((d-1)n+c\bigr)/d, & d\nmid n.\end{cases}
\]
\end{lemma}

\begin{proof}
Both identities are the definition of \(C\) and \(R\), split by the
decidable predicate \(d\mid n\).
\end{proof}

\begin{lemma}[Residues]\label{Lem:Residues}
Let \(0<d\) and \(d\nmid n\). Then
\[
 d\mid (d-1)n+c \iff n\equiv c \pmod d ,
 \qquad
 (d-1)n+c\not\equiv c \pmod d .
\]
In particular \(n\equiv c\pmod d\) implies \(d\mid(d-1)n+c\).
\end{lemma}

\begin{proof}
In \(\mathbb{Z}/d\), the affine step satisfies
\((d-1)n+c\equiv -n+c\pmod d\), so \(d\mid(d-1)n+c\) holds exactly when
\(-n+c\equiv0\), that is \(n\equiv c\pmod d\). For the second assertion,
\((d-1)n+c\equiv c\) means \(-n\equiv0\), that is \(d\mid n\), which fails
by hypothesis; the residue class \(c\) is therefore never attained by the
affine step at a state that is not a multiple of \(d\).
\end{proof}

\begin{remark}\label{Rem:Inaccessible}
Lemma~\ref{Lem:Residues} is the reason accessibility is not closed under
\(C\): at an inaccessible state the affine step is again inaccessible, since
it is congruent to \(c\) only at states that are already divisible by \(d\).
\end{remark}

\section{Contraction, the core, and traps}

\begin{lemma}[Free step increases]\label{Lem:Free}
If \(d\nmid n\) and \(c\ge1\), then \(n<(d-1)n+c\).
\end{lemma}

\begin{proof}
\(d\ge2\) gives \((d-1)n\ge n\), and \(c\ge1>0\).
\end{proof}

\begin{lemma}[Division decreases]\label{Lem:DivDec}
If \(d\mid n\) and \(0<n\), then \(n/d<n\).
\end{lemma}

\begin{proof}
Write \(n=d\cdot(n/d)\) with \(n/d\ge1\) and apply \(d\ge2\).
\end{proof}

\begin{theorem}[Strict contraction above the core]\label{Thm:Contract}
Let \(d\ge2\), \(c\ge1\) and \(c<n\). Then \(R(n)<n\), and if \(d\mid n\)
then \(C(n)<n\).
\end{theorem}

\begin{proof}
If \(d\mid n\), then \(R(n)=C(n)=n/d<n\) by
Lemma~\ref{Lem:DivDec}. Otherwise \(d\nmid n\), so \(R(n)=((d-1)n+c)/d\).
Since \(d\nmid n\) implies \(n\ge c+1\) or \(n<c\): in the first case
\((d-1)n+c<dn\) because \(c<n\); in the second case \(n\le c-1<c<n\), a
contradiction with \(c<n\). So \((d-1)n+c<dn\) always holds under \(c<n\),
and dividing by \(d\) gives \(R(n)<n\).
\end{proof}

\begin{theorem}[Core invariance]\label{Thm:Core}
Let \(d\ge2\), \(c\ge1\) and \(1\le n\le c\). Then \(1\le R(n)\le c\), and
for every \(k\) the iterate \(R^{k}(n)\) satisfies
\(1\le R^{k}(n)\le c\).
\end{theorem}

\begin{proof}
If \(d\mid n\) then \(1\le n/d\le n\le c\). Otherwise
\(R(n)=((d-1)n+c)/d\). For the upper bound, \(n\le c\) gives
\(((d-1)n+c)/d\le((d-1)c+c)/d=c\). For the lower bound, \(n\ge1\) and
\(c\ge1\) give \((d-1)n+c\ge d-1+1=d\), hence \(R(n)\ge1\). The iterate
statement follows by induction on \(k\), since the bounds \(1\le\cdot\le c\)
are preserved at every step.
\end{proof}

\begin{theorem}[Traps]\label{Thm:Trap}
Let \(d\ge2\), \(c\ge1\), and let \(n\) be inaccessible. Then for every
\(k\) the iterate \(C^{k}(n)\) is inaccessible, strictly increasing,
\(C^{k}(n)<C^{k+1}(n)\), and \(n+k\le C^{k}(n)\); consequently for every
bound there is \(j\) with \(C^{j}(n)>\text{bound}\).
\end{theorem}

\begin{proof}
Lemma~\ref{Lem:Residues} shows that the successor of an inaccessible state is
again inaccessible, so by induction on \(k\) every \(C^{k}(n)\) is
inaccessible. The free branch therefore applies at every step, and
Lemma~\ref{Lem:Free} gives \(C^{k}(n)<C^{k+1}(n)\) for all \(k\). For the
linear bound, induct on \(k\): the case \(k=0\) is trivial, and the step uses
\(C^{k+1}(n)-C^{k}(n)\ge1\). Given any bound \(b\), choose \(j=b\) if
\(b>n\) and \(j=\max(b,n)+1\) otherwise; then \(C^{j}(n)\ge n+j>b\).
\end{proof}

\section{Realization and cycle expansion}

\begin{theorem}[Realizing return steps]\label{Thm:Realize}
Let \(d\ge2\), \(1\le x\), and let \(R^{k}(x)\) be accessible for every
\(k<p\). Then there is \(j\) with \(p\le j\le 2p\) and
\(C^{j}(x)=R^{p}(x)\). Moreover one instruction realizes the return step at a
multiple of \(d\), and two instructions realize it at a positive state with
\(d\nmid x\) and \(x\equiv c\pmod d\).
\end{theorem}

\begin{proof}
At a multiple of \(d\), \(C(x)=x/d=R(x)\). At a state with \(d\nmid x\) and
\(x\equiv c\pmod d\), Lemma~\ref{Lem:Residues} gives \(d\mid(d-1)x+c\), so
\(C^{2}(x)=C((d-1)x+c)=((d-1)x+c)/d=R(x)\). Induction on \(p\) over the
accessible return prefix therefore realizes each return step by one or two
instructions, and the accumulated count \(j\) satisfies \(p\le j\le2p\).
\end{proof}

\begin{corollary}[Cycle expansion]\label{Cor:Expand}
Let \(d\ge2\), \(1\le x\), \(p>0\), let \(R^{k}(x)\) be accessible for every
\(k<p\), and let \(R^{p}(x)=x\). Then there is \(j\) with \(0<j\le2p\) and
\(C^{j}(x)=x\).
\end{corollary}

\begin{proof}
Apply Theorem~\ref{Thm:Realize} and substitute \(R^{p}(x)=x\).
\end{proof}

\section{Reachability, periodicity, and the dichotomy}

\begin{theorem}[Reaching the core]\label{Thm:Reach}
Let \(d\ge2\), \(c\ge1\) and \(1\le n\). Then \(R^{k}(n)\le c\) for some
\(k\le n\); by Theorem~\ref{Thm:Core} that value lies in
\(\{1,\dots,c\}\).
\end{theorem}

\begin{proof}
If \(n\le c\) take \(k=0\). Otherwise \(c<n\), and
Theorem~\ref{Thm:Contract} gives \(R(n)<n\), so by strong induction on \(n\)
the state \(R(n)\) reaches the core within \(n-1\) return steps; since every
intermediate value is at most its predecessor, the total is at most \(n\).
\end{proof}

\begin{theorem}[Periodicity]\label{Thm:Periodic}
Let \(d\ge2\), \(c\ge1\) and \(1\le n\). Then there are \(k\) and \(p>0\)
with
\[
 R^{k+s+p}(n)=R^{k+s}(n)\quad\text{for every }s,
\]
and \(R^{k}(n)\in\{1,\dots,c\}\).
\end{theorem}

\begin{proof}
By Theorem~\ref{Thm:Reach} some \(R^{k}(n)\le c\) with
\(R^{k}(n)\ge1\). The map \(R\) sends the finite set \(\{1,\dots,c\}\) into
itself by Theorem~\ref{Thm:Core}, so the \(c+1\) states
\(R^{k}(n),\dots,R^{k+c}(n)\) lie in a set of \(c\) elements and two of them
coincide. If \(R^{k+i}(n)=R^{k+j}(n)\) with \(i<j\), then \(p=j-i>0\) and
\(R^{k+s+p}(n)=R^{k+s}(n)\) for every \(s\), by applying the same number of
return steps to the two sides.
\end{proof}

\begin{theorem}[Trap-or-core dichotomy]\label{Thm:Dichotomy}
Let \(d\ge2\), \(c\ge1\) and \(1\le n\). Then either some iterate
\(C^{k}(n)\) is inaccessible, and then \(C^{k+t}(n)\) is inaccessible and
passes every bound as \(t\to\infty\), or every iterate \(C^{k}(n)\) is
accessible.
\end{theorem}

\begin{proof}
Both alternatives are decided by the first iterate that is inaccessible, if
one exists; Theorem~\ref{Thm:Trap} then supplies growth and unboundedness.
\end{proof}

\section{The corner case \(c=d+1\)}

\begin{lemma}[Corner residues]\label{Lem:CornerResidue}
If \(d\ge2\), then \((d+1)\bmod d=1\).
\end{lemma}

\begin{proof}
\(d+1=1\cdot d+1\) with \(1<d\).
\end{proof}

\begin{theorem}[Corner core states]\label{Thm:CornerCore}
Let \(d\ge3\) and \(1\le x\le d+1\). The state \(x\) is accessible if and
only if \(x\in\{1,d,d+1\}\). Consequently every \(2\le x<d\) is inaccessible.
\end{theorem}

\begin{proof}
Accessibility means \(d\mid x\) or \(x\equiv d+1\equiv1\pmod d\). In the first
case \(x\in\{d\}\) because multiples of \(d\) in \(\{1,\dots,d+1\}\) are \(d\)
only. In the second case \(x\bmod d=1\) forces \(x=1\) or \(x=d+1\) in
\(\{1,\dots,d+1\}\); the intermediate states \(2\le x<d\) satisfy
\(x\not\equiv0\) and \(x\not\equiv1\pmod d\) because \(1<x<d\). Conversely \(d\) is
a multiple of \(d\), \(1\equiv d+1\pmod d\), and \(d+1\) equals \(d+1\) modulo
\(d\), so each of \(1\), \(d\) and \(d+1\) is accessible.
\end{proof}

\begin{theorem}[Corner values and traps]\label{Thm:CornerValues}
Let \(d\ge2\). Then \(R(1)=2\), \(R(d)=1\), and \(R(d+1)=d+1\), and
\(d+1\) is the only accessible fixed point of \(R\) in
\(\{1,\dots,d+1\}\). Let \(d\ge3\). Then every \(x\) with \(2\le x<d\)
is inaccessible.
\end{theorem}

\begin{proof}
\(d\nmid1\) requires \(d\ge2\) and \(1\equiv1\equiv c\pmod d\), and
\(((d-1)+d+1)/d=2d/d=2\). Further \(d\mid d\) and \(R(d)=1\). Finally
\(d\nmid d+1\), so
\(R(d+1)=((d-1)(d+1)+d+1)/d=d(d+1)/d=d+1\). The trap range is
Theorem~\ref{Thm:CornerCore}. For uniqueness, an accessible \(x\le d+1\) is
\(1\), \(d\) or \(d+1\) by Theorem~\ref{Thm:CornerCore} when \(d\ge3\), and
at \(d=2\) the accessible states \(1,2,3\) satisfy \(R(1)=2\), \(R(2)=1\) and
\(R(3)=3\). In both cases the first two are not fixed points since \(R(1)=2>1\)
and \(R(d)=1<d\), while \(R(d+1)=d+1\).
\end{proof}

\begin{theorem}[The corner cycle]\label{Thm:CornerCycle}
Let \(d\ge3\). Then \(C(d+1)=d(d+1)\) and \(C(d(d+1))=d+1\). Consequently
\[
 C^{2m}(d+1)=d+1,\quad C^{2m+1}(d+1)=d(d+1),\quad
 C^{2m}(d(d+1))=d(d+1),\quad C^{2m+1}(d(d+1))=d+1
\]
for every \(m\), and every iterate of either cycle state is one of the two
states.
\end{theorem}

\begin{proof}
\(d\nmid d+1\) gives \(C(d+1)=(d-1)(d+1)+d+1=d(d+1)\), and
\(d\mid d(d+1)\) gives \(C(d(d+1))=d+1\). The iterate formulas follow by
induction on \(m\), and the last claim is immediate from them.
\end{proof}

\begin{theorem}[Corner classification]\label{Thm:CornerClass}
Let \(d\ge3\) and \(1\le x\). Then either some iterate \(C^{k}(x)\) is
inaccessible, and then \(C^{k+t}(x)\) passes every bound as
\(t\to\infty\), or \(C^{k}(x)\in\{d+1,d(d+1)\}\) for some \(k\). In the
second case the orbit reaches the cycle.
\end{theorem}

\begin{proof}
If some iterate is inaccessible, Theorem~\ref{Thm:Trap} gives the divergence.
Assume then that every iterate is accessible. As long as the states exceed
\(d\) they are not in the trap range of Theorem~\ref{Thm:CornerValues}, are
positive, and Theorem~\ref{Thm:Contract} applies to every return step: once
the state is larger than \(c=d+1\), one return step strictly decreases it.
Theorem~\ref{Thm:Reach} therefore supplies \(k\) with
\(R^{k}(x)\le d+1\), and \(R^{k}(x)\) is accessible, hence
\(R^{k}(x)\in\{1,d,d+1\}\) by Theorem~\ref{Thm:CornerCore}. Since every
iterate is accessible, the trap alternatives \(1\) and \(d\) are excluded:
\(R(1)=2\) and \(R(d)=1\) lead to the inaccessible state \(2\) (for \(d\ge3\)),
contradicting accessibility of all iterates. Hence \(R^{k}(x)=d+1\), and
Theorem~\ref{Thm:Realize} realizes the return iterates by instructions of
\(C\), giving \(C^{j}(x)=d+1\) for some \(j\), which is on the cycle
\(d+1,d(d+1)\) by Theorem~\ref{Thm:CornerCycle}.
\end{proof}

\begin{theorem}[Corner periodic points]\label{Thm:CornerPeriodic}
Let \(d\ge3\), \(1\le x\), and \(k>0\). Then
\[
 C^{k}(x)=x \iff \bigl(x=d+1\ \text{or}\ x=d(d+1)\bigr)\ \text{and}\ 2\mid k .
\]
\end{theorem}

\begin{proof}
If \(x\in\{d+1,d(d+1)\}\) and \(2\mid k\), Theorem~\ref{Thm:CornerCycle}
gives \(C^{k}(x)=x\). Conversely let \(C^{k}(x)=x\). If some iterate of \(x\)
is inaccessible, then Theorem~\ref{Thm:Trap} makes the orbit unbounded after
that iterate, contradicting periodicity: if \(C^{j}(x)\) is inaccessible with
\(j<k\), then \(C^{k}(x)=C^{k-j}(C^{j}(x))\ge C^{j}(x)+(k-j)>C^{j}(x)\), and
repeating the argument over multiples \(tk\) of the period contradicts
boundedness. Hence every iterate is accessible, and
Theorem~\ref{Thm:CornerClass} yields \(j\) with \(C^{j}(x)\in\{d+1,d(d+1)\}\).
Iterating the period identity gives \(C^{k t}(x)=x\) for every \(t\); choosing
\(t\) with \(kt\ge j\) and applying Theorem~\ref{Thm:CornerCycle} to the
cycle state \(C^{j}(x)\) over the remaining \(kt-j\) instructions gives
\(C^{kt-j}(C^{j}(x))\in\{d+1,d(d+1)\}\), so \(x\in\{d+1,d(d+1)\}\). Finally
\(C^{k}(x)=x\) with \(x\) on a cycle of length two forces \(2\mid k\) by
Theorem~\ref{Thm:CornerCycle}.
\end{proof}

\begin{remark}[Periodic-point count]\label{Rem:CornerCount}
The number of positive fixed points of \(C^{k}\) is \(2\) for even \(k\) and
\(0\) for odd \(k\), so a single raw cycle of length two carries the
periodic-point zeta function \(1/(1-z^{2})\) on the positive states. The state
\(0\) is fixed as well, so a count over all nonnegative fixed points would
differ at odd \(k\).
\end{remark}

\section{Controls and excluded cases}

\begin{itemize}[leftmargin=2em]
\item The state \(0\) is fixed by every \(C\) and lies in the domain, so it is
periodic without reaching the corner cycle; positivity is a genuine
hypothesis of Theorems~\ref{Thm:CornerClass} and \ref{Thm:CornerPeriodic}.
\item Theorem~\ref{Thm:Reach} carries \(c\ge1\) in the formal statement because
its proof uses positivity of the core; the return map itself still reaches the
core at \(c=0\), since it strictly decreases on positive states. At \(d=2\),
\(c=0\) the raw map satisfies \(C(1)=1\cdot1+0=1\), so that single affine
instruction does not increase.
\item Theorem~\ref{Thm:CornerValues} needs no \(d\ge3\) for its values and its
uniqueness, which hold at \(d\ge2\), and its trap range is empty at \(d=2\).
The classification needs \(d\ge3\): at \(d=2\), \(c=3\) the state \(2\) is
accessible (\(2\equiv0\pmod2\)) and the map carries the further cycle
\(1\mapsto4\mapsto2\mapsto1\) beside \(3\mapsto6\mapsto3\), so
Theorem~\ref{Thm:CornerClass} would be false with two raw cycles.
\item The claim that the divergent set has natural density one is not proved
here. The certificate counts, below \(300000\), the seeds whose corner orbit
reaches the cycle: \(1977\), \(63\), \(57\), \(22\), \(31\), \(16\), \(15\) and
\(10\) for \(d=3,4,5,6,7,9,12,20\). Finite counting is not a density
statement, and the density claim remains open.
\end{itemize}

\section{Formal correspondence}

\begin{table}[htbp]
\centering
\begin{tabular}{ll}
\toprule
Statement & Lean declaration \\
\midrule
Lemma~\ref{Lem:AffineBranch} & \texttt{CoreMapOfDvd}, \texttt{CoreMapOfNotDvd}, \texttt{CoreReturnOfDvd}, \texttt{CoreReturnOfNotDvd} \\
Lemma~\ref{Lem:Residues} & \texttt{CoreMapDivisibleIff}, \texttt{CoreMapResidueIff}, \texttt{CoreDomainDivisible}, \texttt{CoreMapCast} \\
Lemma~\ref{Lem:Free} & \texttt{CoreMapFreeIncrease} \\
Lemma~\ref{Lem:DivDec} & \texttt{CoreMapDivisibleDecrease} \\
Theorem~\ref{Thm:Contract} & \texttt{CoreReturnDecrease} \\
Theorem~\ref{Thm:Core} & \texttt{CoreReturnCore}, \texttt{CoreReturnIterateCore} \\
Theorem~\ref{Thm:Trap} & \texttt{CoreMapNotDomain}, \texttt{CoreTrapNotDomain}, \texttt{CoreTrapStrictMono}, \texttt{CoreTrapGrowth}, \texttt{CoreTrapUnbounded} \\
Theorem~\ref{Thm:Realize} & \texttt{CoreReturnRealizeDvd}, \texttt{CoreReturnRealizeMod}, \texttt{CoreReturnRawRealization} \\
Corollary~\ref{Cor:Expand} & \texttt{CoreReturnCycleRaw} \\
Theorem~\ref{Thm:Reach} & \texttt{CoreReturnReachesCore}, \texttt{CoreReturnReachesCoreBounded} \\
Theorem~\ref{Thm:Periodic} & \texttt{CoreReturnEventuallyPeriodic} \\
Theorem~\ref{Thm:Dichotomy} & \texttt{CoreMapDichotomy} \\
Lemma~\ref{Lem:CornerResidue} & \texttt{CoreCornerModSelf} \\
Theorem~\ref{Thm:CornerCore} & \texttt{CoreCornerCoreStates}, \texttt{CoreCornerCoreStatesIff} \\
Theorem~\ref{Thm:CornerValues} & \texttt{CoreReturnCornerOne}, \texttt{CoreReturnCornerD}, \texttt{CoreReturnCornerSelf}, \texttt{CoreMapCornerTrap}, \texttt{CoreReturnCornerFixedUnique} \\
Theorem~\ref{Thm:CornerCycle} & \texttt{CoreMapCornerCycleSteps}, \texttt{CoreMapCornerCycle}, \texttt{CoreMapCornerCycleBack}, \texttt{CoreCornerCycleIterateEven}, \texttt{CoreCornerCycleIterateOdd}, \texttt{CoreCornerCycleIterateBack}, \texttt{CoreCornerCycleIterateBackOdd}, \texttt{CoreCornerCycleInvariant} \\
Theorem~\ref{Thm:CornerClass} & \texttt{CoreCornerClassification}, \texttt{CoreCornerDomainForeverHitsCycle}, \texttt{CoreCornerIterateRealized} \\
Theorem~\ref{Thm:CornerPeriodic} & \texttt{CoreCornerPeriodicPointCycle}, \texttt{CoreCornerPeriodicIff} \\
Remark~\ref{Rem:Inaccessible} & \texttt{CoreMapNotDomain} \\
Remark~\ref{Rem:CornerCount} & \texttt{CoreCornerPeriodicIff} \\
\bottomrule
\end{tabular}
\caption{Lean declarations carrying each statement.}
\end{table}

\end{document}
